
\subsection{Limitations}
% 1. Model assumptions
% DSHARP dust composition
% grain-size distribution
% spherical vs. irregular grain structure
% filling factor/porosity
% the scattering model itself

% 2. Limited range of grain structures
% - You tested a range of porosities for spherical grains, but Glitterin was used for compact irregular grains.


% 3. Polarization mechanisms aren't uniquely identified
% - morphology alone can't prove which mechanism is operating.

% 4. Optical depth interpretation
% You've been using optical depth as an explanation for the spatial change. But unless directly calculated a detailed optical-depth map, not certain


% 5. Observational limitations / uncertainties
% finite angular resolution
% sensitivity
% polarization detection thresholds
% limited number of wavelengths
% uncertainties in polarization fraction

\subsection{Future Work}
% More wavelengths, especially to better resolve transition.
% Full radiative-transfer modelling with realistic disk structure/optical depth.
% Broader grain composition/size distribution/shape/porosity parameter space.
% More irregular/porous grain modelling with Glitterin.
% Higher-resolution polarization observations.
% Apply method to more Class I disks to determine whether early grain growth is common.